\chapter{Dépôt et reproduction des résultats}
\label{ann:depot}

\section{Organisation du dépôt}

Le dépôt sépare strictement ce qui est saisi de ce qui est calculé
(figure~\ref{fig:depot}).

\begin{figure}[htbp]
  \centering
  % Organisation du dépôt. Arbre vectoriel plutôt que listing ASCII : la
% couleur y distingue ce qui est versionné de ce qui est régénéré.
\begin{tikzpicture}[font=\sffamily\scriptsize,
    dossier/.style={draw=navy, fill=navy!6, text=navy, rounded corners=2pt,
                    line width=0.6pt, minimum height=5.4mm, inner xsep=5pt,
                    font=\sffamily\scriptsize\bfseries, anchor=west},
    fichier/.style={draw=rule, fill=white, text=ink, rounded corners=2pt,
                    line width=0.5pt, minimum height=5.4mm, inner xsep=5pt,
                    font=\ttfamily\scriptsize, anchor=west},
    genere/.style={fichier, draw=slate!45, fill=mist, text=slate},
    note/.style={font=\sffamily\scriptsize, text=slate, anchor=west},
    branche/.style={draw=rule, line width=0.5pt}]

\def\ind{0.55}   % retrait par niveau
\def\pas{0.62}   % pas vertical

% #1 niveau, #2 ligne, #3 style, #4 libellé, #5 commentaire
\newcommand{\entree}[5]{%
  \node[#3] (n#2) at ({#1*\ind},{-#2*\pas}) {#4};
  \node[note] at ({5.6},{-#2*\pas}) {#5};}

\entree{0}{0}{dossier}{sovereign-legal-rag/}{}
\entree{1}{1}{dossier}{data/}{}
\entree{2}{2}{fichier}{raw/}{PDF sources, jamais modifiés}
\entree{2}{3}{genere}{processed/corpus\_chunks.jsonl}{régénéré par le parseur}
\entree{2}{4}{genere}{chroma/}{base vectorielle, régénérée}
\entree{1}{5}{dossier}{src/}{}
\entree{2}{6}{fichier}{parser.py}{PDF $\rightarrow$ corpus structuré}
\entree{2}{7}{fichier}{database.py}{corpus $\rightarrow$ index dense}
\entree{2}{8}{fichier}{retriever.py}{dense + BM25 + fusion RRF}
\entree{2}{9}{fichier}{generator.py}{prompt ancré, appel Ollama}
\entree{2}{10}{fichier}{app.py}{abstention et contrôle des citations}
\entree{1}{11}{dossier}{eval/}{}
\entree{2}{12}{fichier}{reference\_qa.jsonl}{jeu de référence, \NbQuestions{} questions}
\entree{2}{13}{fichier}{run\_eval.py}{Recall@k et MRR, sans LLM}
\entree{2}{14}{fichier}{run\_answer\_eval.py}{critères portant sur la réponse}
\entree{2}{15}{fichier}{run\_measures.py}{latence et taxonomie}
\entree{2}{16}{genere}{mesures-performance.json}{source des chiffres du rapport}
\entree{1}{17}{dossier}{rapport/}{}
\entree{2}{18}{fichier}{build\_data.py}{mesures $\rightarrow$ macros et tables}
\entree{1}{19}{fichier}{web/ et serve.py}{interface de démonstration}
\entree{1}{20}{fichier}{JOURNAL.md}{carnet de bord quotidien}

% ------------------------------------------------------------------ branches
% Verticale du niveau 1, des enfants directs de la racine.
\draw[branche] ({0.5*\ind},{-0.34}) -- ({0.5*\ind},{-20*\pas});
\foreach \r in {1,5,11,17,19,20}
  \draw[branche] ({0.5*\ind},{-\r*\pas}) -- ({1*\ind-0.03},{-\r*\pas});

% Verticales de niveau 2, une par dossier parent.
\foreach \parent/\premier/\dernier in {1/2/4, 5/6/10, 11/12/16, 17/18/18} {
  \draw[branche] ({1.5*\ind},{-\parent*\pas-0.28}) --
                 ({1.5*\ind},{-\dernier*\pas});
  \foreach \r in {\premier,...,\dernier}
    \draw[branche] ({1.5*\ind},{-\r*\pas}) -- ({2*\ind-0.03},{-\r*\pas});
}

% -------------------------------------------------------------------- légende
\node[fichier, minimum width=0pt, text width=0pt, minimum height=3.4mm]
  (lg1) at (0,{-21.4*\pas}) {};
\node[note, right=1.5mm of lg1] (lt1) {versionné};
\node[genere, minimum width=0pt, text width=0pt, minimum height=3.4mm,
      right=5mm of lt1] (lg2) {};
\node[note, right=1.5mm of lg2] {régénéré (absent du dépôt)};

\end{tikzpicture}

  \caption[Organisation du dépôt de code]{Organisation du dépôt. Les éléments grisés ne sont pas
  versionnés : ils se reconstruisent depuis les sources, ce qui est la condition pour que la
  reproduction annoncée ci-dessous ait un sens.}
  \label{fig:depot}
\end{figure}

\section{Reproduction des résultats}

L'ensemble des chiffres de ce rapport se régénère par les commandes suivantes, exécutées depuis la
racine du dépôt. Les deux premières reconstruisent les données ; les trois suivantes produisent
les mesures ; la dernière convertit ces mesures en macros et coordonnées de figures.

\begin{lstlisting}[language=bash,caption={Commandes de reproduction},label={lst:repro}]
# 1. Construction du corpus a partir du PDF officiel
python src/parser.py

# 2. Indexation (telecharge le modele d'embeddings au premier lancement)
python src/database.py

# 3. Metriques de recherche : Recall@k et MRR, sans appel au LLM
python eval/run_eval.py

# 4. Evaluation des reponses generees (necessite Ollama en fonctionnement)
python eval/run_answer_eval.py

# 5. Mesures de latence et taxonomie des resultats
python eval/run_measures.py

# 6. Conversion des mesures en donnees de rapport
python rapport/build_data.py
\end{lstlisting}

\section{Chaîne des chiffres du rapport}

Le script \module{rapport/build\_data.py} lit deux fichiers en entrée :

\begin{itemize}
  \item \module{eval/mesures-performance.json} --- toutes les mesures d'exécution ;
  \item \module{data/processed/corpus\_chunks.jsonl} --- le corpus structuré.
\end{itemize}

\noindent Il en produit deux sortes de fichiers :

\begin{itemize}
  \item \module{Figures/data/chiffres.tex}, qui définit une macro LaTeX pour chaque nombre cité
        dans le texte --- \verb|\TauxCitation|, \verb|\LatTotale|, \verb|\NbArticles|, etc. ;
  \item \module{Figures/data/fig-*.tex}, qui contiennent les coordonnées lues par les figures.
\end{itemize}

Aucun nombre n'est donc écrit en dur dans les chapitres. Une nouvelle exécution du banc suivie
d'une recompilation met à jour simultanément le texte et les figures, ce qui rend impossible la
divergence habituelle entre une phrase et le graphique censé l'illustrer.

\section{Prérequis d'exécution}

Les prérequis diffèrent selon l'étape (tableau~\ref{tab:prerequis}).

\begin{table}[htbp]
  \centering
  \caption{Prérequis pour rejouer la chaîne complète}
  \label{tab:prerequis}
  \small
  \begin{tabularx}{\textwidth}{@{}>{\raggedright\arraybackslash}p{3.6cm} R@{}}
    \toprule
    \sffamily\bfseries\color{navy}Élément &
    \sffamily\bfseries\color{navy}Remarque \\
    \midrule
    Python 3.11 ou supérieur & Dépendances listées dans \module{requirements.txt} \\
    Ollama & Nécessaire aux étapes 4 et 5 uniquement ; modèle \ModeleLLM{} à télécharger au
             préalable \\
    BGE-M3 & Téléchargé automatiquement au premier lancement de l'étape 2 \\
    PDF source & Présent dans \module{data/raw/}, jamais modifié par la chaîne \\
    \bottomrule
  \end{tabularx}
\end{table}

Les étapes 1 à 3 ne requièrent aucun modèle de génération : la qualité de la recherche peut donc
être reproduite sans Ollama, ce qui était l'un des objectifs de la séparation entre recherche et
génération.
